% ===================================================================
%  اللوحة رقم ٨ — الحصاد. — الصيد، قطاف العنب، صيد السمك.
%  Traduction 2026 d'apres texto/io, controlee sur texto/fr, texto/en
%  et texto/es. Consigne : voir texto/ar/10-lawha-01.tex.
%
%  L'appel de note est dans le TITRE de la scene — « الحصاد (*) » — et
%  non dans un alinea : c'est ainsi que l'ido le compose.
% ===================================================================

%%K t08-noto-1 noto
\VUnotes{143.2mm}{%
\textit{(*) لأن هذه الكلمة غير دقيقة: أهو حصاد كتّان؟ أم قطاف عنب؟ أم
جني تفاح؟ ولعل الأولى أن يُستعمل «} \VUgras{mesono} \textit{» المقترح في}
Mondo \textit{الرابع عشر، ص ٢٤٢. لكن هذا الجذر الجديد لم يعتمده المجمع
حتى الآن، وهو ينتظر المناقشة وفق قواعد حركة الإيدو المعتادة.}%
}

%%K t08-apar-1 apar
\VUsaut{5.0mm}
\VUtitre{15.0pt}{60}{اللوحة رقم ٨}
\VUsaut{3.0mm}
\VUcentre{13.2pt}{90}{\textit{المشهد الأول.}}
\VUsaut{3.0mm}
\VUpk{10.2pt}{40}{الحصاد (*).}
\VUsaut{4.0mm}
\VUpk{10.2pt}{40}{وجه الريف.}
\VUsaut{3.0mm}

%%K t08-c1-01-1 p
1. --- مع \VUgras{الصيف} غيّر الريف وجهه مرة أخرى. فقد أنبت البذر الذي
أُودع الأخدود؛ وخرجت من الأرض نبتة خضراء صغيرة حملت سنبلتها. وتكوّنت
الحبة ثم نضجت، ولم يبق الآن إلا جمع المحصول.

%%K t08-c1-02-1 p
2. --- وقد تفتّحت \VUgras{أزهار البر}. فـ\VUgras{زهر الذرة} \textsuperscript{(1)}
و\VUgras{شقائق النعمان} \textsuperscript{(2)} و\VUgras{الكشتبانية}
\textsuperscript{(3)} تضع على لون حقول القمح المتساوي تباينَ \VUgras{أزهارها}
النضرة المبهج.

%%K t08-c1-03-1 p
3. --- وقد أورقت أشجار الفاكهة؛ ونضج ثمرها. و\VUgras{شجرة الكرز}
\textsuperscript{(4)} الصغيرة مثقلة بـ\VUgras{كرز} \textsuperscript{(5)} جميل يقطفه
الأطفال ويأكلونه في لذة.

%%K t08-c1-04-1 p
4. --- وفي وسع القطيع الآن أن يقضي النهار كله في الخلاء.

%%K t08-c1-04-2 p
فقد كبرت \VUgras{الحملان} \textsuperscript{(6)} وصارت \VUgras{نعاجًا}
\textsuperscript{(7)} جميلة و\VUgras{كباشًا} \textsuperscript{(8)} جميلة ذات أصواف
كثيفة. وهي ترعى في هدوء \VUgras{عشب} \VUgras{المرعى} \textsuperscript{(9)}
المرصّع بـ\VUgras{الأقحوان}.

%%K t08-c1-04-3 p
وعمّا قريب تُجزّ هذه الحيوانات ليؤخذ \VUgras{صوفها}، وهو الذي يُصنع منه
جوخ ثيابنا.

%%K t08-tit-2 sub
\VUsaut{5.0mm}
\VUpk{10.2pt}{40}{راعي البقر.}
\VUsaut{3.4mm}

%%K t08-c1-05-1 p
5. --- ولا يبدو أن \VUgras{راعي البقر} \textsuperscript{(10)} يعيرها كبير
انتباه. فهمّه الوحيد العاصفة المارّة في الأفق، أما \VUgras{راعي الغنم}
\textsuperscript{(13)} فيبدو أشد منه لا مبالاة وهو يرفع \VUgras{قِربته}
\textsuperscript{(14)} إلى شفتيه.

%%K t08-c1-06-1 p
6. --- والحرّ خانق؛ فـ\VUgras{الكلب} \textsuperscript{(15)} يأتي هو أيضًا
ليتبرّد؛ يلعق ماء \VUgras{الجدول} \textsuperscript{(16)} العذب فيفزع
\VUgras{ضفدعة} \textsuperscript{(17)} صغيرة مسكينة كانت رابضة في العشب على
الضفة. فتثب إلى الماء الصافي حيث تتفتح \VUgras{زنابق الماء}
\textsuperscript{(18)}.

%%K t08-c1-07-1 p
7. --- و\VUgras{ذعرة} \textsuperscript{(19)} تستحم قرب \VUgras{النبع}
\textsuperscript{(20)} الصغير الذي ينبثق بين صخرتين ويختفي تحت
\VUgras{الشوك} \textsuperscript{(21)} و\VUgras{أجمات الزعرور} \textsuperscript{(22)}.

%%K t08-tit-3 sub
\VUsaut{5.0mm}
\VUpk{10.2pt}{40}{الحصاد.}
\VUsaut{3.4mm}

%%K t08-c1-08-1 p
8. --- وموسم حصاد القمح موسم مسلٍّ جدًا لأطفال الريف. يقلبون
\VUgras{شقلبات} \textsuperscript{(24)} مرحة على \VUgras{أكوام القش}
\textsuperscript{(23)}، ويعينون \VUgras{الحصّادين} \textsuperscript{(26)}؛ وبينما
يحصد هؤلاء \VUgras{حقل القمح} \textsuperscript{(27)} بـ\VUgras{مناجلهم}
\textsuperscript{(25)} المسنونة جيدًا على \VUgras{المِسنّ} \textsuperscript{(30)}،
يجمع الأطفال القمح ويربطونه \VUgras{حِزَمًا} \textsuperscript{(31)} أو
\VUgras{رُزَمًا صغيرة} \textsuperscript{(32)}.

%%K t08-c1-09-1 p
9. --- وأحيانًا يُؤذن لأكبرهم بأن يقطع بـ\VUgras{منجل} \textsuperscript{(12)}
\VUgras{السيقان الذهبية} \textsuperscript{(28)} التي تحمل \VUgras{السنابل}
\textsuperscript{(29)} الثقيلة الملأى بالحب؛ أما الصغار فيرعون
\VUgras{الماشية} \textsuperscript{(33)} وهي ترعى في \VUgras{المرج}
\textsuperscript{(34)} في ظل \VUgras{الزيزفونة} \textsuperscript{(35)} الكبيرة،
ويصطادون في أثناء ذلك \VUgras{الفراشات} الجميلة بـ\VUgras{شبكتهم}
\textsuperscript{(36)}.

%%K t08-c1-09-2 p
بل إن بعضهم يتسلق \VUgras{العربة} \textsuperscript{(37)} ليرتّب الحزم التي
تُرفع إليه على \VUgras{مِذراة} \textsuperscript{(38)}.

%%K t08-c1-10-1 p
10. --- وفي \VUgras{السهل} \textsuperscript{(39)} كله، حيث تطير
\VUgras{القبّرات} \textsuperscript{(41)}، حركة كبيرة. فالجميع مشغولون بالحصاد.
لكن، بينما يظل الفقراء يستعملون الأدوات القديمة --- \VUgras{المناجل
والمحاش} و\VUgras{المدارس} \textsuperscript{(40)} --- يحصد ملّاك الأرض الأغنياء
قمحهم أو \VUgras{شعيرهم} \textsuperscript{(43)} بـ\VUgras{آلة حصاد}
\textsuperscript{(42)}. ثم تُقام في المكان نفسه \VUgras{أكداس} \textsuperscript{(44)}
كبيرة، دون حاجة إلى نقل الحزم إلى الحظيرة.

%%K t08-c1-11-1 p
11. --- ثم يُؤتى بـ\VUgras{آلة الدرس} \textsuperscript{(47)}، يديرها
\VUgras{محرّك متنقل} \textsuperscript{(45)} عبر \VUgras{سير نقل} \textsuperscript{(46)}،
وهكذا يُفصل الحب عن \VUgras{القش} دون مغادرة الحقل الذي زُرع فيه.

%%K t08-tit-4 sub
\VUsaut{5.0mm}
\VUpk{10.2pt}{40}{العاصفة.}
\VUsaut{3.4mm}

%%K t08-c1-12-1 p
12. --- وكثيرًا ما تنفجر \VUgras{عواصف} \textsuperscript{(53)} عنيفة في موسم
الحصاد. فيُسمع أولًا قصف \VUgras{الرعد} من بعيد؛ وتهبّ \VUgras{ريح غربية}
\textsuperscript{(54)} قوية تُحني أكبر أشجار \VUgras{الغابة} \textsuperscript{(49)}
حتى تئنّ. وتقتحم \VUgras{السحب} \textsuperscript{(56)} السوداء السماء؛ ويشقّها
تعرّج \VUgras{البرق} \textsuperscript{(55)} الخاطف. ويبدأ \VUgras{المطر}، بل
\VUgras{البَرَد} أحيانًا، بالنزول فيُلقي المحصول الجاهز للحصاد أرضًا.

%%K t08-c1-13-1 p
13. --- وكثيرًا ما تشعل الصاعقة بناءً. ففي العام الماضي مثلًا احترق سطح
\VUgras{الطاحونة الهوائية} \textsuperscript{(50)} حتى صار رمادًا، وتحطّم
\VUgras{جناحان} \textsuperscript{(51)} من أجنحتها.

%%K t08-c1-14-1 p
14. --- وبعد ساعات يعود الهدوء. ويظهر في الأفق \VUgras{قوس قزح}
\textsuperscript{(52)} باهر، وتستأنف الطبيعة الحياة التي انقطعت حينًا. ويعود
أهل الريف، وقد احتموا في \VUgras{الأكواخ} \textsuperscript{(48)}، إلى العمل
ويشرعون في شجاعة يصلحون ما أصابهم من ضرر.

%%K t08-tit-5 sub
\VUsaut{6.0mm}
\VUcentre{13.2pt}{90}{\textit{المشهد الثاني.}}
\VUsaut{3.6mm}

%%K t08-tit-6 sub
\VUpk{10.2pt}{40}{الصيد. --- قطاف العنب.}
\VUsaut{1.4mm}
\VUpk{10.2pt}{40}{صيد السمك.}
\VUsaut{4.0mm}

%%K t08-c2-01-1 p
1. --- في أواخر \VUgras{أغسطس} أو أواسط \VUgras{سبتمبر}، على اختلاف
الأقاليم، يُفتح موسم الصيد. فما إن تُخرج أشعة الشمس الأولى
\VUgras{السحالي} \textsuperscript{(2)} من جحورها حتى ينطلق \VUgras{الصياد}
\textsuperscript{(5)} مع \VUgras{كلابه} \textsuperscript{(4)}.

%%K t08-c2-02-1 p
2. --- وقد لبس \VUgras{بذلة الصيد} \textsuperscript{(9)} المتينة من الجوخ
الثخين، و\VUgras{طماقًا} من جلد تمكّنه من السير في العشب المبلول حيث تختبئ
\VUgras{الضفادع} \textsuperscript{(1)}؛ وأخذ \VUgras{بندقيته} \textsuperscript{(6)}
و\VUgras{مخلاته} \textsuperscript{(10)}، ومضى.

%%K t08-c2-03-1 p
3. --- ويقوده حماس الصيد إلى وسط كرم يجري فيه القطاف على أشدّه. فأولاد
الكرّام، وهم يرعون المعز عادة على جانب الطريق، يتركونها اليوم ترعى حيث
شاءت، وبعد أن ربطوا إلى وتد \VUgras{تيسًا} \textsuperscript{(3)} عجوزًا كان
سيغتنم حريته حتمًا ليفرّ، أخذوا هم أيضًا نصيبهم من اللهو. فواحد يأخذ
عنقودًا من \VUgras{سلة العنب} \textsuperscript{(12)} ويتسلى بعصر \VUgras{العصير}
\textsuperscript{(14)} في \VUgras{كأس} \textsuperscript{(13)} ليصنع سُلافة (خمرًا غير
مخمّرة). وآخر يتوّج نفسه بـ\VUgras{إكليل من الورق} \textsuperscript{(15)} جدله
لتوّه. وقربهم تأتي \VUgras{العصافير} \textsuperscript{(16)} الجريئة إلى المعصرة
نفسها لتسرق العنب.

%%K t08-c2-04-1 p
4. --- وبينما يلهو الأطفال يعمل الرجال. فـ\VUgras{القطّافون}
\textsuperscript{(18)} ينزلون بالعنب إلى القبو في \VUgras{سلال ظهر}
\textsuperscript{(17)}؛ و\VUgras{البرميلي} \textsuperscript{(19)} يصلح
\VUgras{البراميل} \textsuperscript{(20)}، ويتفقّد سلامة \VUgras{الألواح}
\textsuperscript{(22)} و\VUgras{القِيعان} \textsuperscript{(21)}، ويبدّل
\VUgras{الأطواق} \textsuperscript{(23)} المكسورة. وهو الذي صنع أيضًا
\VUgras{الأحواض} \textsuperscript{(25)} الكبيرة التي يُصبّ فيها \VUgras{العنب}
\textsuperscript{(26)} ليختمر.

%%K t08-c2-05-1 p
5. --- فإذا انتهى التخمّر سُحبت الخمر ووُضع الباقي على \VUgras{المعصرة}
\textsuperscript{(28)}؛ وبشدّ \VUgras{اللولب} \textsuperscript{(29)} الكبير شيئًا
فشيئًا يُعصر العصير فيجري في \VUgras{حوض} \textsuperscript{(32)}، فتُستخرج
\VUgras{خمر} \textsuperscript{(31)} أخرى.

%%K t08-tit-8 sub
\VUsaut{6.0mm}
\VUpk{10.2pt}{40}{الجعة وعصير التفاح.}
\VUsaut{3.4mm}

%%K t08-c2-06-1 p
6. --- وعلى جدار \VUgras{القبو} \textsuperscript{(27)} يتسلق غصن من
\VUgras{الجنجل} \textsuperscript{(30)}. وهو هنا نبات للزينة فحسب، لكنه يُزرع في
بعض البلدان، حيث الكرمة نادرة جدًا، لصنع \VUgras{الجعة}.

%%K t08-c2-07-1 p
7. --- و\VUgras{شجرة التفاح} \textsuperscript{(33)} الصغيرة مثقلة بتفاح سيعطي
\VUgras{عصير تفاح} ممتازًا متى نضج. وفي \VUgras{قفصهما} \textsuperscript{(34)}
ينظر \VUgras{الكناري} \textsuperscript{(35)} و\VUgras{الحسّون} \textsuperscript{(36)}
بعين حسود إلى العصافير التي تمرح في حرية.

%%K t08-c2-08-1 p
8. --- وعلى امتداد البصر، على سفح التلال، تقوم صفوف من \VUgras{أوتاد
الكرم} \textsuperscript{(40)} تسند \VUgras{جذوع الكرم} \textsuperscript{(39)} البديعة
المثقلة بـ\VUgras{العناقيد}.

%%K t08-c2-08-2 p
وطوال السنة عمل \VUgras{الكرّام} \textsuperscript{(42)} على رعاية \VUgras{كرمه}
\textsuperscript{(38)}، وها هو يُجزى على تعبه بمحصول طيب. و\VUgras{القطّافات}
\textsuperscript{(41)} يملأن الأحواض الموضوعة على العربة التي تجرّها
\VUgras{البغال} \textsuperscript{(43)}، والجميع مبتهجون.

%%K t08-tit-9 sub
\VUsaut{5.0mm}
\VUpk{10.2pt}{40}{صيد السمك.}
\VUsaut{3.4mm}

%%K t08-c2-09-1 p
9. --- وكذلك \VUgras{الصيادون بالصنارة} \textsuperscript{(44)}، وقد جاءوا مع أول
ضوء ليستقروا على حافة الماء بـ\VUgras{قصباتهم} \textsuperscript{(45)}.

%%K t08-c2-09-2 p
فـ\VUgras{السمك} \textsuperscript{(47)} يبتلع الشِّص ما إن يُلقوا \VUgras{خيطهم}
\textsuperscript{(46)} في الماء، وما يكادون يجدّدون \VUgras{الطُّعم} حتى تكون
ضحية جديدة تتخبط في طرف الخيط.

%%K t08-c2-09-3 p
ومع ذلك فـ\VUgras{الصياد بالشبكة} \textsuperscript{(48)}، الذي يُلقي شبكته من
\VUgras{قاربه} \textsuperscript{(49)} إلى وسط \VUgras{النهر} \textsuperscript{(50)}،
يصيد أكثر منهم بكثير.

%%K t08-c2-10-1 p
10. --- والخريف موسم النزهات الطيبة: على الأقدام، أو على \VUgras{ظهر
الخيل} \textsuperscript{(52)}، أو على \VUgras{الدراجة} \textsuperscript{(53)}، أو
بـ\VUgras{السيارة} \textsuperscript{(54)}. و\VUgras{الطرق} \textsuperscript{(51)}
نظيفة، والناس يغتنمون آخر الأيام الصافية، فقد أخذت \VUgras{السنونوات}
\textsuperscript{(56)} تتجمع على السطوح لتطير إلى البلاد الدافئة. وأوراق
\VUgras{شجرة الجوز} \textsuperscript{(57)} العتيقة تصفرّ، و\VUgras{الجوز} يتساقط،
و\VUgras{الأشجار} الباسقة القائمة \VUgras{أيكةً} \textsuperscript{(58)} على
\VUgras{المرتفع} \textsuperscript{(59)} ستتعرّى عن قريب.

%%K t08-c2-11-1 p
11. --- و\VUgras{الشمس} \textsuperscript{(60)} تشرق متأخرة، والنهار يقصر، ويبدأ
\VUgras{الصباح} يُحسّ فيه برد لاذع بعض الشيء، وقد أخذت أسراب
\VUgras{الدُّجّ} \textsuperscript{(61)} تغادر مناخنا غير المضياف.
\VUsaut{6.0mm}
\VUfilet{22mm}
